\documentclass[11pt]{exam}

% Usepackages si setari pentru ecuatii
\usepackage[utf8x]{inputenc}
\usepackage{titlesec}

\everymath{\displaystyle}

\usepackage[margin=1in]{geometry}
\usepackage{amsmath,amssymb}
\usepackage{multicol}
\usepackage{graphicx}

% Gestiunea formatului header-ului
\newcommand{\class}{Olimpiada de Matematică, clasa a X-a }
\newcommand{\theschool}{Liceul Teoretic A.I.Cuza}
\newcommand{\examnum}{Teză semestrul I}
\newcommand{\term}{Anul școlar 2023-2024}
\newcommand{\examdate}{Ianuarie 2024}
\newcommand{\timelimit}{90 de minute}

\firstpageheader{}{}{}
\runningheader{\examnum, disciplina \class}{}{\examdate}
\runningheadrule

% Aranjamentul punctelor
\renewcommand{\points}{p}
% Plaseaza punctele in partea dreapta
\pointsdroppedatright
\pointsinrightmargin

% Micsorare font pentru sectiuni
\titleformat*{\section}{\normalsize\bfseries}
\titleformat*{\subsection}{\normalsize\bfseries}

\begin{document}

\noindent
\begin{tabular*}{\textwidth}{l @{\extracolsep{\fill}} r @{\extracolsep{3pt}} l}
& \textbf{Numele :} & \makebox[2in]{\hrulefill}\\
\textbf{\class} &&\\
\textbf{\examnum} & \textbf{Clasa:} & \makebox[2in]{\hrulefill}\\
\textbf{\term} &&\\
\textbf{\examdate} &&\\
%& \textbf{Nota:} & \makebox[2in]{\hrulefill}
\end{tabular*}

\vspace{10 mm}

%\rule[2ex]{\textwidth}{2pt}

{\setlength{\parindent}{1 mm}
Acest test conține \numquestions\ întrebări
%Numărul total de puncte este \numpoints. 
Pe foaia de examen scrieti redactarile complete.
}

{\setlength{\parindent}{2 mm}
\textbf{Timp de lucru: \timelimit} 
}

\vspace{1mm}

\begin{questions}


\begingradingrange{section1}


\question[7] Fie $z_1, z_2, z_3\in \mathbb{C}$ cu $|z_1| =|z_2|= |z_3| =1$ si $z_1 +z_2 +z_3 =1$. Atunci calculati $z_1^{2023} +z_2^{2023}+z_3^{2023}$.

\question[7] Fie $z_A = 2+5i, z_B= 8-i, z_C= -2+i$
\begin{parts}
\part[3] Aflati afixul centrului cercului circumscris $\Delta ABC$.
\part[4] Aflati afixul punctului $D$, unde $AD$ este inaltimea dusa din $A$ in $\Delta ABC, D\in BC.$
\end{parts}


\endgradingrange{section1}

\setcounter{question}{0}

\end{questions}

\end{document}

